\documentclass[11pt]{article}
\usepackage[margin=1in]{geometry}
\usepackage{amsmath,amssymb,amsthm}
\usepackage[hidelinks]{hyperref}
\usepackage{xurl}
\setlength{\emergencystretch}{2em}
\newtheorem{theorem}{Theorem}
\newtheorem{lemma}[theorem]{Lemma}
\newtheorem{corollary}[theorem]{Corollary}
\newcommand{\rk}{\operatorname{rank}}
\title{A signed rank-one obstruction at every meridian parameter}
\author{Internal research working proof}
\date{17 November 2026}
\begin{document}\maketitle

For $13a_{373}$, the ordinary $U=1$ augmentation rank over $\mathbb C$
is three, whereas its meridional rank and bridge number are four.
The upper bound follows from a signed rank-one matrix obstruction valid
over every field and at every nonzero meridian determinant. It does not
assume a preserved form or specialize the meridian parameter.

\begin{theorem}\label{thm:generic}
Let $k$ be any field. Every irreducible representation of the knot group
of $13a_{373}$ whose meridian matrices have displacement rank one has
dimension at most three. In particular, the knot has no irreducible
four-dimensional KCH representation at any meridian parameter
$\mu\in k^*\setminus\{1\}$.
No preserved bilinear or Hermitian form is assumed.
\end{theorem}

The argument below is accompanied by exact algebra controls. Its diagram input is
the retained literal PD for $13a_{373}$, together with the named-diagram
binding from the seven-sheet package. The signed crossing extraction in
the checker follows the oriented literal planar diagram and uses the
following Wirtinger convention:
\[
 M_b=M_o^\varepsilon M_aM_o^{-\varepsilon},
 \qquad\varepsilon=-\text{crossing sign}
\]
when $a,b$ are the oriented incoming and outgoing under-arcs. This is
also the convention in Spherogram's explicit knot-group construction.
Reversing
the two under-arcs reverses $\varepsilon$.

Each signed tuple below is $(o,a,b,\varepsilon)$, with over-arc $o$
and under-arc relation $M_b=M_o^\varepsilon M_aM_o^{-\varepsilon}$.
The seed arcs are $1,2,6,8$. Their signed propagation trace is
\[
\begin{gathered}
(2,6,7,+),\ (7,2,3,+),\ (3,1,0,+),\ (1,3,4,-),\\
(8,4,5,-),\ (5,8,9,-),\ (6,9,10,+),\\
(10,0,12,+),\ (4,10,11,-).
\end{gathered}
\]
The unused signed crossings are
\[
 (0,1,2,-),\quad(12,5,6,-),\quad
 (11,7,8,+),\quad(9,11,12,+).
\]
Only the first three of these are needed below.

Write $d=1-\mu$ and
\[
 M_i=I-u_if_i,\qquad f_i(u_i)=d,\qquad
 M_i^{-1}=I+\mu^{-1}u_if_i.
\]
This holds for every invertible rank-one displacement: the determinant
is $1-f_i(u_i)=\mu\ne0$, the same for all conjugate meridians.
For exponent $+$ the pair propagates as
\[
 u_b=M_ou_a,\qquad f_b=f_aM_o^{-1},
\]
and for exponent $-$ as
\[
 u_b=M_o^{-1}u_a,\qquad f_b=f_aM_o.
\]
Equality of two nonzero rank-one displacements identifies their pairs
up to $(u,f)\mapsto(cu,c^{-1}f)$ with $c\ne0$.

The span of the four seed roots is invariant under the four generating
meridians. Their covectors' common kernel is fixed pointwise.
Irreducibility therefore bounds the dimension by four; in dimension four
the seed roots and seed covectors are both bases. It suffices to exclude
that case. Take the root basis to be $A,B,C,D$ on arcs $1,2,6,8$, with
covector basis $F_A,F_B,F_C,F_D$.

\section*{The first signed clasp}
Put
\[
a=F_B(C),\quad b=F_C(B),\quad r=F_A(B),\quad R=F_B(A),
\quad x=F_A(C),\quad X=F_C(A).
\]
Define
\[
k=1+ab/\mu,\qquad S=kR+aX,\qquad T=kr-bx/\mu.
\]
The first three positive propagation steps give
\[
\begin{aligned}
u_3&=kB-(b/\mu)C,& f_3&=kF_B+aF_C,\\
u_0&=A-Su_3,&f_0&=F_A+(T/\mu)f_3.
\end{aligned}
\]
Let $H=f_0(A)=d+ST/\mu$. The negative closing crossing
$(0,1,2,-)$ says
\[
 M_0^{-1}A=\kappa B,\qquad F_AM_0=\kappa^{-1}F_B
\]
for some $\kappa\ne0$. The root equation has coefficient vector
\[
 (1+H/\mu)A-(HS/\mu)kB+(HSb/\mu^2)C.
\]
Thus $H=-\mu$, so $ST=-\mu$, and $Sk=\kappa$, $Sb=0$.
In particular $S,T\ne0$, whence $b=0$ and $k=1$.
Also $F_A(u_0)=d-ST=1$, so the covector equation becomes
\[
 -(T/\mu)(F_B+aF_C)=\kappa^{-1}F_B.
\]
Independence forces $a=0$. Consequently $rR=-\mu$.
Rescale $B\mapsto RB$ and $F_B\mapsto R^{-1}F_B$; then
$F_B(A)=1$ and $F_A(B)=-\mu$.

The pairs on arcs $0,4,6,7,8$ now have the form
\[
 (U,F),\quad(-U,-F),\quad(V,G),\quad(V,G),\quad(W,J),
\]
where $U=A-B$, $F=F_A-F_B$, $V=C$, $G=F_C$, $W=D$, $J=F_D$.
Both triples are linearly independent, and their corresponding
evaluations are $F(U)=G(V)=J(W)=d$.

\section*{The three-direction core}
Allow six independent directed pairings
\[
x=F(V),\quad X=G(U),\quad y=F(W),\quad Y=J(U),\quad
z=G(W),\quad Z=J(V).
\]
No symmetry is imposed. Set
\[
\begin{gathered}
h=1+yY/\mu,\qquad v=yX+hz,\qquad C=hZ-Yx/\mu,\\
w=y(d+h)-xv,\qquad Q=Y(h-d/\mu)+CX/\mu.
\end{gathered}
\]
The remaining signed propagation steps give the following coefficient
vectors in the independent triples $(U,V,W)$ and $(F,G,J)$:
\[
\begin{array}{c|c|c}
i&u_i&f_i\\\hline
5&(-1,0,-Y/\mu)&(-1,0,y)\\
9&(y,0,h)&(-Y/\mu,0,h)\\
10&(y,-v,h)&(-Y/\mu,C/\mu,h)\\
11&(y+w/\mu,-v,h)&(-Y/\mu-Q,C/\mu,h)\\
12&(1-Qy,Qv,-Qh)&(1-wY/\mu^2,wC/\mu^2,wh/\mu).
\end{array}
\]
The exact checker verifies each row using the signed conjugation formula
and verifies every corresponding evaluation $f_i(u_i)=d$.

The positive crossing $(11,7,8,+)$ supplies a nonzero scalar $\beta$:
\[
 M_{11}V=\beta W,\qquad GM_{11}^{-1}=\beta^{-1}J.
\]
Put $p=f_{11}(V)$, $q=G(u_{11})$. Coefficient comparison gives
\[
\begin{array}{lll}
p(y+w/\mu)=0,&1+pv=0,&-ph=\beta,\\
(q/\mu)(-Y/\mu-Q)=0,&1+qC/\mu^2=0,&qh/\mu=\beta^{-1}.
\end{array}
\]
These equations force $p,q,h\ne0$, and hence
\[
w=-\mu y,\quad Q=-Y/\mu,\quad v=\beta^{-1}h,
\quad C=-\mu\beta h.
\]
Substitution simplifies the arc12 pair to
\[
u_{12}=-hu_5-(Yv/\mu)V,\qquad
f_{12}=-hf_5-(yC/\mu)G.
\]
Finally the negative crossing $(12,5,6,-)$ supplies $\alpha\ne0$ with
\[
 M_{12}^{-1}u_5=\alpha V,\qquad
 f_5M_{12}=\alpha^{-1}G.
\]
Write $r=f_{12}(u_5)$ and $s=f_5(u_{12})$.
The independent pairs $(u_5,V)$ and $(f_5,G)$ yield
\[
1-rh/\mu=0,\quad-rYv/\mu^2=\alpha,\qquad
1+sh=0,\quad syC/\mu=\alpha^{-1}.
\]
It follows that
\[
 \alpha=-Y\beta^{-1}/\mu,\qquad\alpha^{-1}=y\beta.
\]
Their product forces $yY=-\mu$, so $h=1+yY/\mu=0$,
contradicting $-ph=\beta\ne0$. This proves the theorem.
The proof uses no division by two or by $1-\mu$.

\section*{Exact algebra and augmentation consequence}
The standalone standard-library script
\path{verify_generic_kch_obstruction.py} checks the raw signed
first loop, the normalization, all five reduced paired propagation
steps, and the closing coefficient rows as rational-function identities.
It also constructs an integer Laurent identity
\[
                    \sum_{i=0}^{10}h_if_i=1
\]
for the eleven clean necessary closure equations above. Only
$\mu,\alpha,\beta$ are inverted; there are ten nonzero multiplier
monomials after cancellation. Thus the reduced contradiction works in
every characteristic. This algebra control does not replace the written
independence and signed-diagram argument.

The checker also verifies a genuinely generic four-dimensional
irreducible $T(4,5)$ control, rejects corrupted algebra and diagram
inputs, and verifies the rank-three lower witness described below.
An optional \texttt{--census} run re-reads the pinned named diagram.
These are exact internal verification controls, not formal verification
or external specialist review.

\begin{corollary}\label{cor:augmentation}
For every field $k$, the ordinary $U=1$ augmentation rank satisfies
$\operatorname{ar}(13a_{373},k)\le3$.
\end{corollary}
\begin{proof}
The augmentation consequence is separate from the representation proof.
Cornwell's augmentation--KCH correspondence
\cite[Theorem1.2 and Definition1.8]{cornwell17} applies
to every field-valued $U=1$ cord augmentation with
$\epsilon(\mu)\ne1$ and gives an irreducible KCH representation whose
dimension is its augmentation rank, with arbitrary nonzero longitude
value. Therefore every such augmentation for $13a_{373}$ has rank at
most three. The definition of augmentation rank excludes $\epsilon(\mu)=1$.
This does not assert absence of higher-rank augmentations at $\mu=1$,
outside that correspondence.
\end{proof}

\begin{corollary}\label{cor:gap}
\[
 \operatorname{ar}(13a_{373},\mathbb C)=3
 <4=\operatorname{mr}(13a_{373})=b(13a_{373}).
\]
\end{corollary}
\begin{proof}
Corollary~\ref{cor:augmentation} supplies the augmentation upper bound.
The explicit reflective cord matrix given by the vectors below and in
\path{reflective_augmentation_inputs.json} has rank three over
$\mathbb Q(\phi)$, where $\phi^2-\phi-1=0$. Its thirteen normal vectors
have squared length four, and their half-dot-product matrix satisfies
$c_{ii}=2$ and $c_{\ell a}+c_{\ell b}=c_{\ell o}c_{oa}$ at every crossing.
The checker replays these equations and the matrix rank exactly.
Cornwell's cord presentation and reflective quotient
\cite[Section3]{cornwell15} identify this with a rank-three $U=1$
augmentation at $\mu=-1$, giving equality over $\mathbb C$.
Independently, the seven-sheet cover certificate gives
$\operatorname{mr}\ge4$. The four-seed Wirtinger colouring displayed
above gives Wirtinger number at most four; the bridge--Wirtinger
theorem \cite[Theorem1.3]{wirt} therefore gives $b\le4$.
Thus $\operatorname{mr}=b=4$; the separate cover proof is in
\path{GENERAL_BRANCHED_COVER_RANK_BOUNDS.tex},
Theorem \texttt{thm:twentyfour}.
The meridional lower bound is not used in Theorem~\ref{thm:generic}.
\end{proof}

\begin{corollary}\label{cor:sums}
Let $K=13a_{373}$, and let $s\ge0$ be an integer. Then
\[
 \operatorname{ar}(K\mathbin{\#}(\mathbin{\#}^{s}3_1),\mathbb C)
 =s+3<s+4
 =\operatorname{mr}(K\mathbin{\#}(\mathbin{\#}^{s}3_1))
 =b(K\mathbin{\#}(\mathbin{\#}^{s}3_1)).
\]
\end{corollary}
\begin{proof}
The rank-three seed augmentation is attained at $\mu=-1$.
For the trefoil, the matrices
\[
 A=\begin{pmatrix}-1&1\\0&1\end{pmatrix},\qquad
 B=\begin{pmatrix}1&0\\1&-1\end{pmatrix}
\]
satisfy $A^2=B^2=I$ and $ABA=BAB$. Their displacement roots span
two dimensions, their covectors have zero common kernel, and the
cross-pairings are nonzero, so this is an irreducible two-dimensional
KCH representation at $\mu=-1$. It gives a rank-two augmentation
there. The trefoil has augmentation rank exactly two because its
bridge number is two.

Cornwell's connected-sum theorem \cite[Theorem1.9, proof in Section4.2]{cornwell17}
constructs an augmentation of rank $d_1+d_2-1$ at the common meridian
parameter of the two inputs. Iteration therefore attains rank $s+3$
at $\mu=-1$. Its upper-additivity conclusion starts from an arbitrary
augmentation of the sum and gives the opposite inequality, with no
restriction on that augmentation's parameter. Hence the augmentation
rank is exactly $s+3$. The independent cover and bridge certificates in
\path{INVOLUTION_COVER_SUM_FORMULAS.tex}, Corollary \texttt{cor:hundredone},
give meridional rank and bridge number $s+4$. They do not use the
augmentation argument.
\end{proof}

\section*{The exact diagram and lower witness}
The following PD uses edge labels $1,\ldots,26$ and crossing rows
numbered from zero. Join entries in zero-based positions $1$ and $3$
of every row, then take connected components of edge labels; these are
the overpass arcs, ordered by their least edge label and numbered from
zero. This fixes the arc numbering used in the signed trace.
The retained orientation starts at corner zero of crossing zero.
The exact planar dart bijection to the independently packaged named
census diagram remains in the input file, rather than being inferred
from the knot label.

The second table gives the thirteen witness vectors $v_i$ indexed by
arc. These row indices are independent of the crossing-row indices in
the first table. Here $\phi^2-\phi-1=0$ and
$c_{ij}=\tfrac12v_i\cdot v_j$.
\[
\begin{array}{c|rrrr}
 i&\multicolumn{4}{c}{\mathrm{PD}_i}\\\hline
0 & 3 & 1 & 4 & 26\\
1 & 1 & 7 & 2 & 6\\
2 & 7 & 3 & 8 & 2\\
3 & 13 & 4 & 14 & 5\\
4 & 5 & 14 & 6 & 15\\
5 & 21 & 9 & 22 & 8\\
6 & 9 & 17 & 10 & 16\\
7 & 17 & 11 & 18 & 10\\
8 & 11 & 25 & 12 & 24\\
9 & 19 & 12 & 20 & 13\\
10 & 15 & 22 & 16 & 23\\
11 & 23 & 18 & 24 & 19\\
12 & 25 & 21 & 26 & 20\\
\end{array}
\qquad
\begin{array}{c|rrrr}
 i&\multicolumn{4}{c}{v_i}\\\hline
0 & -1 & 1 & 1 & 1\\
1 & 2 & 0 & 0 & 0\\
2 & -1 & -1 & -1 & -1\\
3 & 1 & 1 & 1 & 1\\
4 & 1 & -1 & -1 & -1\\
5 & \phi-1 & -\phi & 1 & 0\\
6 & \phi-1 & -\phi & 1 & 0\\
7 & 1-\phi & \phi & -1 & 0\\
8 & 1 & -\phi & 0 & 1-\phi\\
9 & 0 & -1 & \phi & \phi-1\\
10 & 1 & -\phi & 0 & 1-\phi\\
11 & \phi-1 & 0 & -\phi & -1\\
12 & 1-\phi & \phi & -1 & 0\\
\end{array}
\]
All $v_i$ have squared length four. Their dot products with
$(0,\phi-2,1-\phi,1)$ vanish, so their span and cord matrix have
rank at most three. The principal cord minor on arcs $0,1,5$ has
\[
                \det(c_{ij})_{i,j\in\{0,1,5\}}=2(1+\phi)\ne0.
\]
For every crossing $(o,a,b)$ and every $\ell$, substitution gives
$c_{\ell a}+c_{\ell b}=c_{\ell o}c_{oa}$; the checker verifies all
$13^2$ equations exactly in $\mathbb Q[\phi]/(\phi^2-\phi-1)$.
Thus the lower cord-rank witness can be checked from the displayed
vectors without the input file. The seven-sheet meridional lower
certificate and connected-sum cover theorem remain the separate
precisely cited results used in the two corollaries.

\section*{The generic positive knot control}
The positive control is the closure of $(\sigma_1\sigma_2\sigma_3)^5$.
Let $D_4$ be the tridiagonal matrix with diagonal $1-\mu$,
superdiagonal $-\mu$, and subdiagonal one. Put
\[
 T_i=I-e_i(\text{row}_iD_4),\qquad
 \det D_4=1-\mu+\mu^2-\mu^3+\mu^4.
\]
The determinant is nonzero over $\mathbb Q(\mu)$, and both adjacent
directed pairings are nonzero. Every nonzero invariant subspace
therefore contains some $e_i$, then every $e_i$, proving irreducibility.
The retained literal fifteen-crossing PD has roots and covectors whose
signed matrix crossing relations all hold over $\mathbb Z[\mu^{\pm1}]$;
all four $T_i$ occur among its arc images. Its four retained seed roots
have determinant one, and its seed covectors have determinant $\det D_4$.
An exact oriented dart map binds that PD to the displayed closed braid.
Thus the same signed diagram logic accepts an actual generic rank-four
knot representation. The checker rejects a deliberately reversed
individual crossing sign in this noninvolutory control.

Strict inequality between augmentation rank and meridional rank was
already established by Cornwell: the retained arXiv:1310.7526v3,
Corollary1.12 and Theorem5.3, gives
$2=\operatorname{ar}(K,\mathbb C)<\operatorname{mr}(K)=3$ for
$K\in\{8_{16},8_{17},10_{91},10_{94}\}$ \cite{cornwell17}.
The result here supplies the named $3<4$ equality and a signed
rank-one obstruction without a preserved form. No claim that this is
the first augmentation-rank gap is made.

\begin{thebibliography}{9}
\bibitem{wirt}
R. Blair, A. Kjuchukova, R. Velazquez and P. Villanueva,
\emph{Wirtinger systems of generators of knot groups},
Communications in Analysis and Geometry 28(2) (2020), 243--262,
\url{https://doi.org/10.4310/CAG.2020.v28.n2.a2}.
\bibitem{cornwell17}
C. R. Cornwell, \emph{KCH representations, augmentations, and
$A$-polynomials}, Journal of Symplectic Geometry \textbf{15} (2017), 983--1017.
Theorem numbering is that of retained
\href{https://arxiv.org/abs/1310.7526v3}{arXiv:1310.7526v3}.
\bibitem{cornwell15}
C. R. Cornwell, \emph{Character varieties of knot complements and
branched double-covers via the cord ring}, retained
\href{https://arxiv.org/abs/1509.04962v1}{arXiv:1509.04962v1}.
\end{thebibliography}

\end{document}
